% fig_mcl.tex --- MCL laser2d tracking error over a closed-loop run (SYNTHETIC data).
% Data: docs/paper/data/mcl_tracking.csv, produced by experiments/mcl_tracking.cpp,
% which drives the real prism_loc_core::ParticleFilter + Laser2DLikelihoodField with
% KLD resampling, exactly as localization_node.cpp wires it. Every number below comes
% from that CSV. Requires in preamble:
%   \usepackage{pgfplots}\pgfplotsset{compat=1.17}\usepgfplotslibrary{groupplots}
\begin{figure}[t]
  \centering
  \footnotesize
  \begin{tikzpicture}
    \begin{groupplot}[
        group style={group size=1 by 2, vertical sep=1.0cm},
        width=\linewidth,
        height=0.42\linewidth,
        xmin=0, xmax=65,
        xtick={0,10,20,30,40,50,60},
        tick label style={font=\footnotesize},
        label style={font=\footnotesize},
        legend style={font=\footnotesize, at={(0.98,0.97)}, anchor=north east,
                      draw=black!40, fill=white, fill opacity=0.85,
                      text opacity=1, row sep=0.5pt, /tikz/every even column/.append style={column sep=6pt}},
        legend columns=3,
        legend cell align=left,
        axis on top,
        grid=major,
        grid style={gray!20},
      ]
      % --- top: position error (3 seeds) ---
      \nextgroupplot[
          ymin=0, ymax=0.42,
          ytick={0,0.1,0.2,0.3,0.4},
          ylabel={Position error (m)},
        ]
      \addplot[black, solid,  line width=0.5pt]
        table [col sep=comma, x=t, y=pos_err_s0] {figures/../data/mcl_tracking.csv};
      \addlegendentry{seed 101}
      \addplot[black, dashed, line width=0.5pt]
        table [col sep=comma, x=t, y=pos_err_s1] {figures/../data/mcl_tracking.csv};
      \addlegendentry{seed 202}
      \addplot[black, dotted, line width=0.7pt]
        table [col sep=comma, x=t, y=pos_err_s2] {figures/../data/mcl_tracking.csv};
      \addlegendentry{seed 303}

      % --- bottom: yaw error (3 seeds) ---
      \nextgroupplot[
          ymin=0, ymax=0.24,
          ytick={0,0.05,0.10,0.15,0.20},
          xlabel={Time (s)},
          ylabel={Yaw error (rad)},
        ]
      \addplot[black, solid,  line width=0.5pt]
        table [col sep=comma, x=t, y=yaw_err_s0] {figures/../data/mcl_tracking.csv};
      \addplot[black, dashed, line width=0.5pt]
        table [col sep=comma, x=t, y=yaw_err_s1] {figures/../data/mcl_tracking.csv};
      \addplot[black, dotted, line width=0.7pt]
        table [col sep=comma, x=t, y=yaw_err_s2] {figures/../data/mcl_tracking.csv};
    \end{groupplot}
  \end{tikzpicture}
  \caption{Monte-Carlo-localization (MCL) tracking accuracy of the laser2d core on
    a \emph{synthetic} occupancy grid. A simulated robot is driven around a closed
    oval loop with semi-axes 3.6\,m and 2.6\,m (7.2\,m$\times$5.2\,m extent)
    for 65\,s at 10\,Hz through a 12\,m$\times$10\,m
    room (two interior pillars break pose symmetry); 180-beam scans are generated by
    grid-march ray casting (\texttt{test\_helpers.hpp}), and the odometry is a
    drifting frame produced by integrating noise-corrupted motion increments under
    the \texttt{laser2d.yaml} $\alpha$ model. Every plotted value comes from the CSV
    written by \texttt{experiments/mcl\_tracking.cpp}, which runs the real
    \texttt{prism\_loc\_core::ParticleFilter} with the \texttt{Laser2DLikelihoodField}
    model and KLD resampling, wired (predict\,$\to$\,correct\,$\to$\,resample, gated by
    \texttt{update\_min\_d}/\texttt{update\_min\_a}) exactly as \texttt{localization\_node.cpp}.
    Three fixed seeds are shown as thin lines. Starting from a broad 1\,m prior the
    filter converges within a few updates and then tracks the ground truth for the
    whole loop: across the three seeds, steady-state ($t\ge10\,$s) mean position error
    is $0.105\,$m (max $0.29\,$m) and mean yaw error is $0.032\,$rad ($1.8^{\circ}$,
    max $0.18\,$rad).}
  \label{fig:mcl}
\end{figure}
